\pagebreak

\section{The Standard Library - stdlib.3pl}

    \index{library!stdlib.3pl}
    This library contains modules, procedures and functions common
    to all platforms.
    
    This library must be incorporated by using a {\bf module} statement -
    \begin{lstlisting}[gobble=4, language=threepl]
       module "stdlib.3pl"
    \end{lstlisting}
    
    This library includes {\bf multlib.3pl} and {\bf divlib.3pl}. These are
    not documented separately.

    \subsection{global structs}
        
%        Global struct {\em bram} is created for use with block RAM
%        extended access procedures {\em rwrite2()} and {\em rread2()} and
%        function {\em rmemoryout2()}.
%
%        \begin{lstlisting}[gobble=8, language=threepl]
%           struct(bram,
%               bwidth, "int",  // block data width (in 1st only)
%               twidth, "int",  // total data width (in 1st only)
%               block,  "->",   // memory block
%               next,   "->"    // next bram struct in linked list
%           );
%        \end{lstlisting}
%        
%        Global struct {\em mem\_ad} is created for use with external memory
%        access.
%
%        \begin{lstlisting}[gobble=8, language=threepl]
%           struct (mem\_ad,
%               address,    "uint:20",
%               data,       "bits:18"
%           );
%        \end{lstlisting}
%        
%        It is used by the modules {\bf extramread} and {\bf extramwrite}.
        
        Global structs {\bf ds16}, {\bf ds16e} and {\bf ds32} are
        used for block protocols and for high speed serial channel data.

        \begin{lstlisting}[gobble=8, language=threepl]
           struct(global.ds16,
               data,   "bits:16",  // data word
               k,      "log"       // data word contains K-character in MS byte
           );
        \end{lstlisting}

        \begin{lstlisting}[gobble=8, language=threepl]
           struct(global.ds16e,
               data,   "bits:16",  // data word
               k,      "log",      // data word contains K-character in MS byte
               error,  "log"       // data word is erroneous
           );
        \end{lstlisting}

        \begin{lstlisting}[gobble=8, language=threepl]
           struct(global.ds32,
               data,   "uint:32",  // data word
               k,      "log"       // data word contains K-character in MS byte
           );
        \end{lstlisting}


    \subsection{library modules}
    
        \subsubsection{array\_to\_queue - serialise an array to a
                                                    queue}\label{sec:lmarraytoqueue}


            {\bf Library: stdlib.3pl}

            {\bf Prototype:} \\
            \vsp
            \begin{lstlisting}[gobble=12, language=threepl]
               global module
               array_to_queue (
                   queue(pout),
                   value(busy, "log")
                   <-
                   value(in),
                   value(enable),
                   value(go, "log")
               );
            \end{lstlisting}

            {\bf Output parameters:} \\
            {\bf pout} is the output queue variable\\
            {\bf busy} is true when the module cannot accept another input

            {\bf Input parameters:} \\
            {\bf in} is the input array.\\
            {\bf enable} is an optional argument controlling the number of
            items in {\bf in} to process.\\
            {\bf go} is asserted to trigger the serialisation.

            {\bf description:} \\
            Parameter {\bf in} is an array. It may be value, static, queue,
            selectvalue or priority mode.
            
            Parameter {\bf enable} is a one dimensional array of "log" and may
            be value, static, or selectvalue mode. It must have the same
            dimension as the first or only dimension of parameter {\bf in}.            
            
            When {\bf go} is true the input array is copied to an internal
            structure which is then shifted out to the output queue {\bf pout}.
            The {\bf enable} array indicates which members of array {\bf in} are
            to be copied and shifted, the lowest index false member of {\bf
            enable} indicating that all members of array {\bf in} from that
            index up are ignored. Thus array {\bf enable} must have a number of
            contiguous entries which are true, starting at index 0, and that
            selects the number of contiguous entries in array {\bf in} which
            will be copied and shifter to the output queues, again starting with
            index 0. {\bf go} may contain queue reads.
            
            The second input argument, for parameter {\bf enable}, may be
            omitted in which case all members of {\bf in} will be shifted and
            copied to the output queue.
            
            The output queue must be the same type as the input
            array members. While the serialisation is in progress output value
            {\bf busy} is true. If input {\bf go} is true when {\bf busy} returns to
            false another conversion will take place. If the input {\bf in} is a queue
            then {\bf go} may be held true.
            
            The serialisation takes a number of clock cycles equal to the number
            of entries in the input array that are enabled plus one and hence
            the input data rate is limited to $\frac{1}{n+1}$ where $n$ is the
            number of enabled entries in the input array.



    
        \subsubsection{crc\_8\_16 - 16 bit CRC from 8 bit data}\label{sec:lmcrc816}


            {\bf Library: stdlib.3pl}

            {\bf Prototype:} \\
            \vsp
            \begin{lstlisting}[gobble=12, language=threepl]
               global module
               crc_8_16 (
                   static(out, "uint:16")
                   <-
                   queue(in, "uint:8"),
                   str(poly, "CRC16")
               ) {}
            \end{lstlisting}

            {\bf Output parameters:} \\
            {\bf out} is the output static variable.

            {\bf Input parameters:} \\
            {\bf in} is the input queue.\\
            {\bf poly} is an optional string argument to select the polynomial
            used.

            {\bf description:} \\
            This module accumulate a 16 bit CRC from 8 bit data. Parameter {\bf
            in} is the data input queue, type "uint:8". The CRC is updated as
            each input datum is read. The optional parameter {\bf poly} selects
            one of two polynomials to use. "CRC16" (default) selects $x^{16} +
            x^{15} + x^{13} + 1$. "CCITT" selects $x^{16} + x^{12} + x^5 + 1$. parameter
            {\bf out} is a "uint:16" static variable in which the CRC value is
            accumulated Input data can be read at clock rate.



    
        \subsubsection{crc\_16\_16 - 16 bit CRC from 8 bit data}\label{sec:lmcrc1616}


            {\bf Library: stdlib.3pl}

            {\bf Prototype:} \\
            \vsp
            \begin{lstlisting}[gobble=12, language=threepl]
               global module
               crc_16_16 (
                   static(out, "uint:16")
                   <-
                   queue(in, "uint:16"),
                   str(poly, "CRC16")
               ) {}
            \end{lstlisting}

            {\bf Output parameters:} \\
            {\bf out} is the output static variable.

            {\bf Input parameters:} \\
            {\bf in} is the input queue.\\
            {\bf poly} is an optional string argument to select the polynomial
            used.

            {\bf description:} \\
            This module accumulate a 16 bit CRC from 16 bit data. Parameter {\bf
            in} is the data input queue, type "uint:16". The CRC is updated as
            each input datum is read. The CRC is computed successively from each
            8-bit half of each input datum but in one clock cycle. The optional
            parameter {\bf poly} selects one of two polynomials to use. "CRC16"
            (default) selects $x^{16} + x^{15} + x^{13} + 1$. "CCITT" selects
            $x^{16} + x^{12} + x^5 + 1$. parameter {\bf out} is a "uint:16"
            static variable in which the CRC value is accumulated Input data can
            be read at clock rate.








    
        \subsubsection{crcblock - block 32 bit data with a CRC}\label{sec:lmcrcblock}


            {\bf Library: stdlib.3pl}

            {\bf Prototype:} \\
            \vsp
            \begin{lstlisting}[gobble=12, language=threepl]
               global module
               crcblock (
                   queue(out, ds16)
                   <-
                   queue(in, ds32),
                   int(max_blk_len, 20),
                   int({sob, eob}),
                   log(fixed)
               ) {}
            \end{lstlisting}

            {\bf Output parameters:} \\
            {\bf out} is the output queue, type ds16.

            {\bf Input parameters:} \\
            {\bf in} is the input queue, type ds32.\\
            {\bf max\_blk\_len} is the maximum variable block size or the fixed
                    block size.\\
            {\bf sob} is the start-of-block K-character to be used.\\
            {\bf eob} is the end-of-block K-character to be used.\\
            {\bf fixed} is an optional logical argument to select fixed
                    length blocks.

            {\bf description:} \\
            This module encloses data in a block containing a cyclic reduncancy
            check word. The 16 bit CRC polynomial is $x^{16} + x^{15} + x^{13} +
            1$. Blocks can be fixed or variable length.

            When input data become available the start-of-block word is written
            to the output queue. The start-of-block word has the most significant
            byte set to the start-of-block K character, the least significant
            byte to zero and the K flag is true. If 'fixed' is false (default),
            incoming data is read until either unavailable or max\_block\_len
            items have been read and the block is then terminated. If 'fixed' is
            true, incoming data is read until max\_block\_len items have been
            read and the block is then terminated. each 32 bit word is written
            to the output as a pair of 16 bit words. A 16 bit CRC is accumulated
            from each 16 bit data word. When the block is terminated the 16 bit
            CRC is written to the output followed by an end-of-block word. The
            end-of-block word has the most significant byte set to the
            end-of-block K character, the least significant byte to zero and the
            K flag is true.

            \begin{tabular}{| c | c | c | l}
            \cline{1-3}
            K & upper & lower & \\
            flag & byte & byte & \\
            \cline{1-3}
            \cline{1-3}
            1 & sob &  0  & start-of-block \\ \cline{1-3}
            0 & \multicolumn{2}{| c |}{data word 1} & first data word \\ \cline{1-3}
            0 & \multicolumn{2}{| c |}{data word 2} & second data word \\ \cline{1-3}
            . & \multicolumn{2}{| c |}{.} &  \\
            . & \multicolumn{2}{| c |}{.} &  \\ \cline{1-3}
            0 & \multicolumn{2}{| c |}{data word n} & last data word \\ \cline{1-3}
            0 & \multicolumn{2}{| c |}{CRC} & CRC \\ \cline{1-3}
            1 & eob &  0  & end-of-block \\ \cline{1-3}
            \end{tabular}








    
        \subsubsection{crcunblock - unblock 32 bit data with a CRC}\label{sec:lmcrcunblock}


            {\bf Library: stdlib.3pl}

            {\bf Prototype:} \\
            \vsp
            \begin{lstlisting}[gobble=12, language=threepl]
               global module
               crcunblock (
                   queue(dout, ds32),
                   static(eout)
                   <-
                   queue(in, ds16),
                   int(max_blk_len, 20),
                   int({sob, eob})
               ) {}
            \end{lstlisting}

            {\bf Output parameters:} \\
            {\bf out} is the output queue, type ds32.\\
            {\bf eout} is an optional error counter static variable.

            {\bf Input parameters:} \\
            {\bf in} is the input queue, type ds16.\\
            {\bf max\_blk\_len} is the maximum variable block size or the fixed
                    block size.\\
            {\bf sob} is the start-of-block K-character to be used.\\
            {\bf eob} is the end-of-block K-character to be used.

            {\bf description:} \\
            This module unwraps data from a block containing a cyclic reduncancy
            check word. The 16 bit CRC polynomial is $x^{16} + x^{15} + x^{13} +
            1$. Blocks can be fixed or variable length. The start-of-block word,
            CRC word and end-of-block word are stripped off the input stream and
            data words are packed into 32 bit words (ds32 struct). A CRC is
            computed from the block data including the transmitted CRC word. If
            the resulting CRC is zero the block is written to the output queue.
            If the CRC is not zero the block is discarded.
            
            When a block CRC error is encountered if the output parameter 'eout'
            has an argument the data length of the discarded block is added to
            the value of that variable. Any words received that do not match the
            block format cause 'eout' to be incremented. The size of the 'eout'
            counter is determined by the argument passed in the call to
            crcunblock().

            If the block length exceeds the maximum block length the block is
            terminated and discarded and if the optional error counter output
            parameter is used, the block length is added to the error count.

            Unpacked data is written to a dual-port circular buffer so that
            completed correct data words can be copied to the output data queue
            in parallel with unpacking of incoming blocks. If the output queue
            becomes unavailable for writing, data words will be discarded.


    
  
        \subsubsection{cread - read combinatorial output memory to a
                                                        queue}\label{sec:lmcread}

            {\bf Library: stdlib.3pl}


            {\bf Prototype:} \\
            \vsp
            \begin{lstlisting}[gobble=12, language=threepl]
               global module
               cread (queue(data) <- rmemory(m), int(port), queue(addr)) {}
            \end{lstlisting}

            {\bf Output parameters:} \\
            {\bf data} is the queue variable to be written to from the memory
            read. It takes its type from the argument.

            {\bf Input parameters:} \\
            {\bf m} is the memory.\\
            {\bf port} is the memory port.\\
            {\bf addr} is the memory address.

            {\bf description:} \\
            Read the combinatorial output memory when the input and output
            queues are available. The program logic must ensure that this
            memory port access is mutually exclusive to all other accesses
            to that port.
            
            This module is just a wrapper around a single statement. The source
            is
            \begin{lstlisting}[gobble=12, language=threepl]
               global module
               cread (queue(data) <- rmemory(m), int(port), queue(addr)) {
                   data <<- cast(cmemoryread(m, port, addr), type(data));
               }
            \end{lstlisting}
                        
            It is provided to complete the set {\bf cread()}, {\bf cwrite()},
            {\bf rread()} and {\bf rwrite()} where only rread() contains
            significant code.
            
            See also - \\
            {\bf cwrite()} \autorefph{sec:lmcwrite}, \\
            {\bf rread()} \autorefph{sec:lmrread}, \\
            {\bf rwrite()} \autorefph{sec:lmrwrite}.

    
        \subsubsection{cwrite - write combinatorial output memory from a
                                                    queue}\label{sec:lmcwrite}

            {\bf Library: stdlib.3pl}

            {\bf Prototype:} \\
            \vsp
            \begin{lstlisting}[gobble=12, language=threepl]
               global module
               cwrite (rmemory(m), int(port), queue({addr, data})) {}
            \end{lstlisting}

            {\bf Output parameters:} \\
            none

            {\bf Input parameters:} \\ 
            {\bf m} is the memory.\\
            {\bf port} is the memory port.\\
            {\bf addr} is the memory address.\\            
            {\bf data} is the data to be written.

            {\bf description:} \\
            Write to the combinatorial output memory when the input queues are
            available. The program logic must ensure that this memory port
            access is mutually exclusive to all other accesses to that port.
            
            This module is just a wrapper around a single statement. The source
            is
            \begin{lstlisting}[gobble=12, language=threepl]
               global module
               cwrite (rmemory(m), int(port), queue(addr), queue(data)) {
                   cmemorywrite(m, port, addr, data);
               }
            \end{lstlisting}
                        
            It is provided to complete the set {\bf cread()}, {\bf cwrite()},
            {\bf rread()} and {\bf rwrite()} where only rread() contains
            significant code.
            
            See also - \\
            {\bf cread()} \autorefph{sec:lmcread}, \\
            {\bf rread()} \autorefph{sec:lmrread}, \\
            {\bf rwrite()} \autorefph{sec:lmrwrite}.





        \subsubsection{demultiplex - demultiplex a source\\
                    demultiplex\_l - demultiplex a source using port bit\\
                    demultiplex\_e - demultiplex a source using port number}\label{sec:lmdemultiplex}

            {\bf Library: stdlib.3pl}

            {\bf Prototypes:} \\
            \vsp
            \begin{lstlisting}[gobble=12, language=threepl]
               global module
               demultiplex.... (varargs <- queue(in)) {}
            \end{lstlisting}

            {\bf Output parameters:} \\
            {\bf varargs} any number of output sinks, an array of pointers
            to sinks or a list of sinks.

            {\bf Input parameters:} \\
            {\bf in} is a single input queue.

            {\bf description:} \\
            These modules are intended to demultiplex a single input queue to a
            number of output queues. Outputs need not be the same type but they must
            be assignable from the input type (or {\bf body} field of the input -
            see below). The input need not contain queue reads in which case it will
            be read on every clock cycle.

            These modules may also be called with a single output argument which
            is an array of pointers to queues. In this case they behave as though
            called with the queues as separate arguments. Similarly the single
            output argument may be of type "list", the list containing pointers
            to queues.

            The output types must be a type assignable from the input data
            type. By 'assignable' is meant that the types need not be identical,
            e.g. the number of bits may differ in which case truncation or
            padding will occur or a "uint" may be assigned to an "int".

            For {\bf demultiplex()} the input can be any type. The input is
            copied to the outputs in turn. If an output is not available the
            next available output in cyclic order is selected. If no outputs
            are available the module waits until any output is available and
            this is written to.

            For {\bf demultiplex\_l()} the input type is a wrapper struct
            containing the datum and an output port identifier. That identifier
            is an array of type "log" in which one member is {\bf true} to
            identify the output port. The struct can be created by calling the
            library global function {\bf multiplex\_struct\_l()}
            \autorefph{sec:lfmultiplexstruct}.

            For {\bf demultiplex\_e()} the input type is a wrapper struct
            containing the datum and an output port identifier. That identifier
            is  a "uint" containing the port number (starting at 0). The struct
            can be created by calling the library global function {\bf
            multiplex\_struct\_e()} \autorefph{sec:lfmultiplexstruct}.

            If the selected output queue is an unmatched parameter then the input is
            still read but the data is discarded. For {\bf demultiplex\_l()} and
            {\bf demultiplex\_e()} the directive "demultiplex\_skip\_null\_outputs"
            can be set to true to ignore any unmatched output parameters, but note
            that if an input is directed to that output port the input will not be
            read.

            These modules read and write in the current clock domain.
            Input and output queues may be asynchronous.
            
            The module {\bf demultiplex()} operates by means of a token placed
            on one output. When the input becomes available the first available
            output downstream of the token is written and the token is moved to
            the output following the one written. There are no clock cycles
            lost, as in a traditional token ring, in waiting for the token to
            move around the ring to an available output, thus the module is
            capable of forwarding an input datum on every clock cycle. The
            implementation uses the Xilinx carry chain to minimise the delay
            through the logic cascade of the token ring. If the carry chain were
            not used the logic delays would severely limit the clock rate at
            which this module could be used. Modules {\bf demultiplex\_l()} and
            {\bf demultiplex\_e()} do not need the token ring as the required
            output is explicitly specified by the input.

            
            See also - \\
            {\bf multiplex()} \autorefph{sec:lmmultiplex} \\
            {\bf multiplex\_l()} \autorefph{sec:lmmultiplex} \\
            {\bf multiplex\_e()} \autorefph{sec:lmmultiplex} \\
            {\bf multiplexp()} \autorefph{sec:lmmultiplex} \\
            {\bf multiplexp\_l()} \autorefph{sec:lmmultiplex} \\
            {\bf multiplexp\_e()} \autorefph{sec:lmmultiplex} \\
            {\bf multiplext()} \autorefph{sec:lmmultiplex} \\
            {\bf multiplext\_l()} \autorefph{sec:lmmultiplex} \\
            {\bf multiplext\_e()} \autorefph{sec:lmmultiplex} \\
            {\bf multiplex\_struct\_l()} \autorefph{sec:lfmultiplexstruct} \\
            {\bf multiplex\_struct\_e()} \autorefph{sec:lfmultiplexstruct}.



        \subsubsection{divide - pipeline divide}\label{sec:lmdivide}

            {\bf Library: stdlib.3pl}

            {\bf Prototype:} \\
            \vsp
            \begin{lstlisting}[gobble=12, language=threepl]
               global module
               divide (queue(quot), queue(rem) <- value(num), value(denom)) {}
            \end{lstlisting}

            {\bf Output parameters:} \\
            {\bf quot} is the quotient.\\
            {\bf rem} is the remainder.

            {\bf Input parameters:} \\
            {\bf num} is the dividend.\\
            {\bf denom} is the divisor.
            
            This module is in divlib.3pl which is included in stdlib.3pl.

            {\bf description:} \\
            This module is a pipeline divider. The numerator and denominator
            arguments may be int, uint, int:, uint:, fixed: or ufixed: where the
            trailing ":" indicates target types. For integer types the quotient
            and remainder outputs are integer, signed if either input is signed,
            and either or both the quotient and remainder outputs may be used.
            For one or both inputs fixed point the quotient output will be fixed
            point, signed if either input is signed, and the remainder output
            is not used.
            
            If neither input is a queue or an expression containing a queue the
            divider will accept input on every clock cycle. Either input may be
            immediate mode but both inputs must not be immediate mode. The
            quotient and remainder will become available at the outputs after a
            latency which is determined by the numerator width. Data in transit
            through the module will continue to the outputs regardless of the
            availability of further input data.

            The remainder queue output argument should be at least as wide as the
            denominator argument otherwise the resulting truncation may lose
            data (if the range is known to be constrained then this is
            acceptable).

            For unsigned integer divide (both numerator and denominator
            arguments unsigned target types or non-negative immediate) the
            result width is the width of the numerator argument.

            For signed integer divide (either or both numerator or denominator
            arguments int target types or negative immediate) the result width
            will be one bit wider than for unsigned divide. The extra bit is
            only required for the single case where both numerator and
            denominator are the maximum negative values (e.g. -8 / -1 gives +8,
            but -8 has a width of 4 bits and +8 as a signed int has a width of 5
            bits). If this case is known not to arise a narrower output queue may
            be used.
            
            For fixed point divide the result width is the sum of the numerator
            and denominator widths and the result binary point offset is
            the numerator offset + denominator width - denominator offset.

            The number of pipeline stages is equal to the number of bits in the
            numerator plus one.

            Divide by zero is not detected and gives an undefined result.





        \subsubsection{fifo - a fifo with static interface}\label{sec:lmfifo}

            {\bf Library: stdlib.3pl}

            {\bf Prototype:} \\
            \vsp
            \begin{lstlisting}[gobble=12, language=threepl]
               global module
               fifo (
                   value(out), value({emptyout, fullout}, "log")
                   <-
                   value(in), value({read, write}, "log")
               ) {}
            \end{lstlisting}

            {\bf Output parameters:} \\
            {\bf out} is the data output for a fifo read.\\
            {\bf emptyout} is true if the fifo is empty.\\
            {\bf fullout} is true if the fifo is full.

            {\bf Input parameters:} \\
            {\bf in} is the data input for a fifo write.\\
            {\bf read} causes a fifo read on a clock cycle when high.\\
            {\bf write} causes a fifo write on a clock cycle when high.

            {\bf description:} \\
            Read or write a fifo. On each clock cycle {\bf emptyout} and {\bf
            fullout} indicate when reads and writes respectively cannot be
            executed. The program must ensure that these conditions are
            checked as the module does not do so! If {\bf read} or {\bf
            write} are high in a clock cycle then a read or write will be
            performed. These can be done on the same clock cycle.





        \subsubsection{multiplex - multiplex sources\\
                    multiplex\_l - multiplex sources giving port bit\\
                    multiplex\_e - multiplex sources giving port number\\
                    multiplexp - multiplex queues with priority\\
                    multiplexp\_l - multiplex queues with priority giving port bit\\
                    multiplexp\_e - multiplex queues with priority giving port number\\
                    multiplext - multiplex queues via a tree\\
                    multiplext\_l - multiplex queues via a tree giving port bit\\
                    multiplext\_e - multiplex queues via a tree giving port number}\label{sec:lmmultiplex}

            {\bf Library: stdlib.3pl}

            {\bf Prototype:} \\
            \vsp
            \begin{lstlisting}[gobble=12, language=threepl]
               global module
               multiplex.... (queue(out) <- varargs) {}
            \end{lstlisting}

            {\bf Output parameters:} \\
            {\bf out} is a single output queue.

            {\bf Input parameters:} \\
            {\bf varargs} any number of input sources, an array of pointers
            to sources or a list of sources.

            {\bf description:} \\
            These modules are intended to multiplex a number of input queues to a
            single output queue. Inputs need not be the same type but they must
            be assignable to the output type (or {\bf body} field of the output
            queue - see below). By 'assignable' is meant that the types need not
            be identical, e.g. the number of bits may differ in which case
            truncation or padding will occur or a "uint" may be assigned to an
            "int".          
            
            When inputs become available simultaneously modules {\bf
            multiplex()}, {\bf multiplex\_l()} and {\bf multiplex\_e()} use a
            pseudo-token ring to give equal treatment to all inputs. Note that
            although these modules were written to multiplex queues (or more
            correctly expressions containing queue reads), the input arguments
            could be other modes, say value or static. Where no queue reads are
            involved an expression is always available. In that case the inputs
            would be sampled and copied to the output queue continuously in a
            cyclic order. In principle the inputs could be a mixture of modes in
            which case they would be read cyclically but a queue would be skipped
            if not available for reading.

            Modules {\bf multiplexp()}, {\bf multiplexp\_l()} and {\bf
            multiplexp\_e()} handle simultaneous inputs by choosing the
            left-most available queue in the argument list. For these modules the
            inputs must contain queue reads.

            Modules {\bf multiplext()}, {\bf multiplext\_l()} and {\bf
            multiplext\_e()} handle simultaneous inputs by means of a balanced tree
            structure.
            
            All nine modules may also be called with a single input argument
            which is an array of pointers to sources. In this case they behave
            as though called with the sources as separate arguments. Similarly
            the single input argument may be of type "list", the list containing
            pointers to sources.
            
            In many applications it is necessary to be able to identify the
            origin of an output datum, i.e. the input argument to the module
            from which the datum was read. To facilitate this the module output
            queue type may be a struct which is a wrapper containing the output
            datum and an input port identifier. The output type for {\bf
            multiplex()} and {\bf multiplexp()} do not use a wrapper type, the
            output type being directly compatible with the input type(s).
            
            The output type for {\bf multiplex\_l()}, {\bf multiplexp\_l()} and
            {\bf multiplext\_l()} is a wrapper with two field, {\bf body}
            containing the datum and {\bf port}  which is an array of type "log"
            in which one member is {\bf true} to identify the input port. The
            struct can be created by calling the library global function {\bf
            multiplex\_struct\_l()} \autorefph{sec:lfmultiplexstruct}.
            
            The output type for {\bf multiplex\_e()}, {\bf multiplexp\_e()} and
            {\bf multiplext\_e()} is a wrapper with two fields, {\bf body}
            containing the datum and {\bf port}  which is a "uint" containing
            the port number (starting at 0). The struct can be created by
            calling the library global function {\bf multiplex\_struct\_e()}
            \autorefph{sec:lfmultiplexstruct}.
            
            In the case of {\bf multiplex()}, {\bf multiplexp()} and {\bf
            multiplext()} the inputs are copied to the output queue when
            available. For {\bf multiplex\_l()}, {\bf multiplexp\_l()} and {\bf
            multiplext\_l()} the selected input is copied to the {\bf body}
            field output queue when available and the associated port field has
            the member of the array indexed by the port number assigned true.
            For {\bf multiplex\_e()}, {\bf multiplexp\_e()} and {\bf
            multiplext\_e()} the same applies except that the port number is
            assigned to the "uint" {\bf port} field.
            
            The modules read and write in the current clock domain. Input
            and output queues may be asynchronous.
            
            The modules {\bf multiplex()}, {\bf multiplex\_l()} and {\bf
            multiplex\_e()} operate by means of a token placed on one input.
            When one or more inputs become available the first input downstream
            of the token is read and the token is moved to the input following
            the one read. There are no clock cycles lost, as in a traditional
            token ring, in waiting for the token to move around the ring to an
            available input, thus the module is capable of forwarding an input
            datum on every clock cycle. The implementation uses the Xilinx carry
            chain to minimise the delay through the logic cascade of the token
            ring. Unfortunately the carry chain must be circular and the link
            back from the top of a linear carry chain to the bottom introduces
            a significant delay. The propagation time around the token ring
            is large enough to severely lower the clock rate that can be used
            as the number of inputs increases.
            
            The modules {\bf multiplexp()}, {\bf multiplexp\_l()} and {\bf
            multiplexp\_e()} use a priority mode variable to select which input
            to read. For FPGAs priority mode variables are implemented using the
            Xilinx carry chain if the number of inputs exceeds the maximum
            number of LUT inputs. This is to minimise the delay through cascaded
            logic.
            
            The modules {\bf multiplext()}, {\bf multiplext\_l()} and {\bf
            multiplext\_e()} use a balanced tree with nodes of degree 2 to
            converge the inputs to the output. Each tree node arbitrates between
            its two inputs based on the number of leaves converging on each
            branch. The disadvantage of this scheme is that arbitration is
            based on the fixed configuration, not on the traffic through
            each branch. For example a four input multiplexer where
            input 4 is idle will arbitrate among the other three unfairly
            giving input 3 half the clock cycles, inputs 1 and 2 sharing the
            other half. The advantage of this multiplexer is that it will
            operate at a high clock rate since the arbitration is local to
            the nodes.
                        
            See also - \\
            {\bf demultiplex()} \autorefph{sec:lmdemultiplex} \\
            {\bf demultiplex\_l()} \autorefph{sec:lmdemultiplex} \\
            {\bf demultiplex\_e()} \autorefph{sec:lmdemultiplex} \\
            {\bf multiplex\_struct\_l()} \autorefph{sec:lfmultiplexstruct} \\
            {\bf multiplex\_struct\_e()} \autorefph{sec:lfmultiplexstruct}.
        




        \subsubsection{multiply - pipelined multiply}\label{sec:lmmultiply}

            {\bf Library: stdlib.3pl}

            {\bf Prototype:} \\
            \vsp
            \begin{lstlisting}[gobble=12, language=threepl]
               global module
               multiply (queue(prod) <- value(m1), value(m2)) {}
            \end{lstlisting}

            {\bf Output parameters:} \\
            {\bf prod} is the product.

            {\bf Input parameters:} \\
            {\bf m1} is the multiplicand.\\
            {\bf m2} is the multiplier.
            
            This module is in multlib.3pl which is included in stdlib.3pl.

            {\bf description:} \\
            This module is a pipeline multiplier for signed or unsigned
            integers. The module is useful for Virtex, Virtex E and Spartan 2
            FPGAs which do not contain hardware multipliers. Later Virtex and
            Spartan3 FPGAs containing hardware multipliers offer no reason to
            use this module.
            
            If neither input is a queue or an expression containing a queue the
            multiplier will accept input on every clock cycle unless the output
            queue becomes unavailable. Either input may be a constant but both
            inputs must not be constant. If either input is signed (a target int
            or a -ve constant) the output queue must be signed (an int, not a
            uint).

            The product will become available at the output after a latency
            which is determined by the input widths. Data in transit through the
            module will continue to the output regardless of the availability of
            further input data. The number of pipeline stages depends on the
            size and types of input arguments as follows -

            \begin{tabular}{ l l l }
            unsigned  & expr x expr     & bits in smallest expr \\
            signed    & expr x expr     & bits in smallest expr + 1 \\
            unsigned  & expr x constant & 1-bits in constant \\
            signed    & expr x constant & 1-bits in constant + 1 \\
            \end{tabular}
            
            where signed means that either input argument is signed. {\bf
            expr} means a target variable or expression, possibly containing
            queues. {\bf constant} means an immediate constant, variable or
            expression. Either input argument may be the constant.
            
            The width of the output queue argument should be at least the sum
            of the widths of the two input arguments. If narrower the result
            will be truncated at the most significant end. If the range of
            the product is known to be constrained then truncation of the
            product is acceptable.
            
            


        \subsubsection{pfifo - a prewrite fifo with static interface}\label{sec:lmpfifo}

            {\bf Library: stdlib.3pl}

            {\bf Prototype:} \\
            \vsp
            \begin{lstlisting}[gobble=12, language=threepl]
               global module
               pfifo (
                   value(out),
                   value({emptyout, fullout}, "log")
                   <-
                   value(in),
                   value({read, prewrite, write}, "log"),
                   value(reset, "log")
               ) {}
            \end{lstlisting}

            {\bf Output parameters:} \\
            {\bf out} is the data output for a fifo read.\\
            {\bf emptyout} is true if the fifo is empty.\\
            {\bf fullout} is true if the fifo is full.

            {\bf Input parameters:} \\
            {\bf in} is the data input for a fifo write.\\
            {\bf read} causes a fifo read on a clock cycle when high.\\
            {\bf prewrite} notifies the fifo that a write will be done
            in the future, reserving a place. When no free places remain the
            full condition will be come true.\\
            {\bf write} causes a fifo write on a clock cycle when high. The
            empty condition takes account of data actually written, not
            notified writes (prewrites).\\
            {\bf reset} is an optional reset to clear the fifo.

            {\bf description:} \\
            This fifo is intended to be used where data is to be passed
            through a fifo but there is a delay between execution of the
            data-producing operation and availability of that data. An
            example of this is external memory reads where the data only
            becomes available a number of cycles after the read command has
            been sent. The resulting data are then to be buffered in this
            fifo. To ensure that the memory reads can be pipelined it is
            necessary to notify the fifo that a write will occur later so
            that its test for space can take that into account. Once the
            memory read has been signaled it cannot be stopped so the fifo
            must have space available for the resulting datum some cycles
            later.
            
            On each clock cycle {\bf emptyout} and {\bf fullout} indicate
            when reads and writes respectively cannot be executed. The
            program must ensure that these conditions are checked as the
            module does not do so! If {\bf read}, {\bf prewrite} or {\bf
            write} are high in a clock cycle then a read, prewrite or write
            will be performed. Any combination of these can all occur on
            one clock cycle.


        \subsubsection{pdelay - delay pipeline data}\label{sec:lmpdelay}

            {\bf Library: stdlib.3pl}

            {\bf Prototype:} \\
            \vsp
            \begin{lstlisting}[gobble=12, language=threepl]
               global module
               pdelay (
                   queue(out)
                   <-
                   queue(in),
                   int(delay),
                   value(reset, "log")
               ) {}
            \end{lstlisting}

            {\bf Output parameters:} \\
            {\bf out} is the data output.

            {\bf Input parameters:} \\
            {\bf in} is the data input.\\
            {\bf delay} is the numnber of data to be delayed.\\
            {\bf reset} is an optional reset.

            {\bf description:} \\
            This module delays data by the number specified by argument
            {\bf delay}. The module will only operate when both input
            and output queues are available. Initially {\em delay} input data
            will be consumed while outputing default values (0, false etc.).
            Following this the output datum will be that input {\em delay}
            data cycles ago. The optional reset can be used to return the
            module to its initial condition.


        \subsubsection{rread - read registered output memory to a
                                                    queue}\label{sec:lmrread}

            {\bf Library: stdlib.3pl}

            {\bf Prototype:} \\
            \vsp
            \begin{lstlisting}[gobble=12, language=threepl]
               global module
               rread (queue(data) <- rmemory(m), int(port), queue(addr)) {}
            \end{lstlisting}

            {\bf Output parameters:} \\
            {\bf data} is the queue variable to be written to from the memory
            read. It takes its type from the argument.

            {\bf Input parameters:} \\
            {\bf m} is the memory.\\
            {\bf port} is the memory port.\\
            {\bf addr} is the memory address.

            {\bf description:} \\
            Read the registered output memory when the input and output
            queues are available. There is a latency of 2 cycles between input
            and output.
            
            The program logic must ensure that this
            memory port access is mutually exclusive to all other accesses
            to that port.
            
            See also - \\
            {\bf cread()} \autorefph{sec:lmcread}, \\
            {\bf cwrite()} \autorefph{sec:lmcwrite}, \\
            {\bf rwrite()} \autorefph{sec:lmrwrite}.


        \subsubsection{rwrite - write registered output memory from a
                                                    queue}\label{sec:lmrwrite}

            {\bf Library: stdlib.3pl}

            {\bf Prototype:} \\
            \vsp
            \begin{lstlisting}[gobble=12, language=threepl]
               global module
               rwrite (rram(m), int(port), queue({addr, data})) {}
            \end{lstlisting}

            {\bf Output parameters:} \\
            none

            {\bf Input parameters:} \\
            {\bf m} is the memory.\\
            {\bf port} is the memory port.\\
            {\bf addr} is the memory address.\\            
            {\bf data} is the data to be written.

            {\bf description:} \\
            Write to the registered output memory when the input queues are
            available. The program logic must ensure that this memory port
            access is mutually exclusive to all other accesses to that port.

            The registered output of the port is also changed to the written
            value.
            
            This module is just a wrapper around a single statement. The source
            is
            \begin{lstlisting}[gobble=12, language=threepl]
               global module
               rwrite (rmemory(m), int(port), value(addr), value(data)) {
                   rmemorywrite(m, port, addr, data);
               }
            \end{lstlisting}
                        
            It is provided to complete the set {\bf cread()}, {\bf cwrite()},
            {\bf rread()} and {\bf rwrite()} where only rread() contains
            significant code.
            
            See also - \\
            {\bf cread()} \autorefph{sec:lmcread}, \\
            {\bf cwrite()} \autorefph{sec:lmcwrite}, \\
            {\bf rread()} \autorefph{sec:lmrread}.


        \subsubsection{sort - generate a parallel Batcher sort network}\label{sec:lmsort}

            {\bf Library: stdlib.3pl}

            {\bf Prototype:} \\
            \vsp
            \begin{lstlisting}[gobble=12, language=threepl]
               global module
               sort (queue(out) <- queue(in)) {}
            \end{lstlisting}

            {\bf Output parameters:} \\
            {\bf out} is the output queue.

            {\bf Input parameters:} \\
            {\bf in} is the input queue.

            {\bf description:} \\
            This module generates a parallel sort network. The input and output
            arguments are queues which are 1-dimensional arrays of type uint or
            int. The output ordering depends on the The input and output array
            dimensions and types must be the same. The number of pipeline
            stages in this module depends on the dimensions of the queues. The
            sort network is generated using Batcher's algorithm. Note that the
            amount of logic used grows rapidly with the dimensionality of the
            queues!



        \subsubsection{sqrt - target fixed point square root}\label{sec:lmsqrt}

            {\bf Library: stdlib.3pl}

            {\bf Prototype:} \\
            \vsp
            \begin{lstlisting}[gobble=12, language=threepl]
               global module
               sqrt (queue(out) <- queue(in)) {}
            \end{lstlisting}

            {\bf Output parameters:} \\
            {\bf out} is the output queue.

            {\bf Input parameters:} \\
            {\bf in} is the input queue.

            {\bf description:} \\
            This module performs a bitwise iterative square root on an unsigned
            integer or unsigned fixed point value. The input and output
            arguments must be queue mode variables. The input and output types
            must be both "uint:" or both "ufixed:". The integer or fractional
            widths of input and output need not match. For integers the output
            width will usually be half the input width. For fixed point the
            output type would usually have the point moved left with respect to
            the input. If a result value exceeds the integer or fractional width
            of the output type, left or right truncation will occur.

            Where the operand is not an exact square the result is truncated at
            the least significant end, i.e. the result is the largest value
            whose square is not larger than the operand.
            
            The module does not process data at clock rate; it processes an
            operand in a fixed time of n/2+2 clock cycles where 'n' is the width
            of the (internal) operand.
            
            For an integer square root 'n' is the width of the input, plus 1
            if the width is odd.
            
            For a fixed point square root 'n' is calculated as follows -
            \begin{lstlisting}[gobble=12, language=threepl]
               int(n, size(in));       // total operand width - initially that of input
               int(ifb, offset(in));   // input fraction part width
               int(iib, n - ifb);      // input integer part width
               int(ioffset);           // input left offset to form internal integer
               int(ofb, offset(out));  // output fraction part width
               int(xoffset);           // extra offset for more fractional precision

               if ((ifb % 2) != 0) {   // if input fractional part width odd --
                   n++;                // increment n to make even
                   ioffset = 1;        // input must now be offset by 1 bit
               }
               xoffset = (2 * ofb) - (ifb + ioffset); // extra precision for output fractional part
               if (xoffset > 0) {      // if extra width for output fractional precision --
                   n += xoffset;       // add to n
                   ioffset += xoffset; // add to input offset
               }
               if ((n % 2) != 0)   // if total width odd --
                   n++;            // increment n to make even
            \end{lstlisting}
            This increments n for either integer or fraction part of odd width
            and also widens the the internal operand to account for extra
            precision in the output fractional part compared to the input.
            
            



        \subsubsection{stretch - generate a long pulse from a transition}\label{sec:lmstretch}

            {\bf Library: stdlib.3pl}

            {\bf Prototype:} \\
            \vsp
            \begin{lstlisting}[gobble=12, language=threepl]
               global module
               stretch (static(out, "log") <- value(in, "log"), int(n)) {}
            \end{lstlisting}

            {\bf Output parameters:} \\
            {\bf out} is the output static variable.

            {\bf Input parameters:} \\
            {\bf in} is the input value.\\
            {\bf n} is the required length of the output pulse.

            {\bf description:} \\
            When input argument {\bf in} goes from false to true generate an
            output {\bf out} which is true for the number of cycles specified
            by argument {\bf n}.


        \subsubsection{timebase - generate decadic pulses}\label{sec:lmtimebase}

            {\bf Library: stdlib.3pl}

            {\bf Prototype:} \\
            \vsp
            \begin{lstlisting}[gobble=12, language=threepl]
               global module
               timebase(
                   static({tb_1us, tb_10us, tb_100us, tb_1ms, tb_10ms, tb_100ms, tb_1s}, "log")
                   <-
               ) {}
            \end{lstlisting}

            {\bf Output parameters:} \\
            {\bf tb\_1us} is a $1\mu S$ interval pulse stream. \\
            {\bf tb\_10us} is  a $10\mu S$ interval pulse stream. \\
            {\bf tb\_100us} is  a $100\mu S$ interval pulse stream. \\
            {\bf tb\_1ms} is  a $1mS$ interval pulse stream. \\
            {\bf tb\_10ms} is  a $10mS$ interval pulse stream. \\
            {\bf tb\_100ms} is  a $100mS$ interval pulse stream. \\
            {\bf tb\_1s} is  a $1S$ interval pulse stream. 

            {\bf Input parameters:} \\
            None.

            {\bf description:} \\
            This module generates pulse streams at decade intervals of $1\mu S$
            up to $1S$. Each pulse is one cycle in the current clock domain.
            Each output pulse is aligned with the pulses from higher frequency
            outputs. The current clock frequency must be $\ge$ 2MHz.
            
            The $1\mu S$ pulses are derived from the current clock by integer
            division and hence if the current clock is not an integer
            multiple of 1MHz the interval will differ from $1\mu S$. Each of
            the other outputs is derived from the next higher frequency
            output by division by 10.


        \subsubsection{to\_pulse - generate a pulse from a
                                                    transition}\label{sec:lmtopulse}

            {\bf Library: stdlib.3pl}

            {\bf Prototype:} \\
            \vsp
            \begin{lstlisting}[gobble=12, language=threepl]
               global module
               to_pulse (static(out, "log") <- value(in, "log")) {}
            \end{lstlisting}

            {\bf Output parameters:} \\
            {\bf out} is the output static variable.

            {\bf Input parameters:} \\
            {\bf in} is the input value.

            {\bf description:} \\
            Generate a pulse on {\bf out} when the input argument {\bf in}
            changes from false to true.



        \subsubsection{tws\_recv - two wire serial receiver}\label{sec:lmtwsrecv}

            {\bf Library: stdlib.3pl}


            {\bf Prototype:} \\
            \vsp
            \begin{lstlisting}[gobble=12, language=threepl]
               tws_recv(
                   queue(out)
                   <-
                   value(sck, "log"), value(sdo, "log"), varargs
               ) {}
            \end{lstlisting}

            {\bf Output parameters:} \\
            {\bf out} is the output queue variable.

            {\bf Input parameters:} \\
            {\bf sck} is the input serial clock. \\
            {\bf sdo} is the input serial data. \\
            {\bf varargs} matches two optional arguments, {\bf silence} and
            {\bf settle}.

            {\bf description:} \\
            This module receives a two-wire serial stream and produces
            a parallel output on a queue. Input {\bf sck} is the input clock
            and input {\bf sdo} is the serial data. The data is captured
            on the positive edges of the clock. The clock is low between
            words.
            
            Optional argument {\bf silence} (first argument matched by varags)
            specifies the minimum time between serial words. If {\bf sck}
            goes high before the specified interval has elapsed since the end
            of the current word, the current word will be discarded. The
            argument may be an immediate integer, in which case it specifies
            the number of clock cycles in the current clock domain, or it
            may be immediate floating point, which will be interpreted as
            a time in second. If the argument is missing or is zero no
            minimum time period between words will be enforced.
            
            Optional argument {\bf settle} (second argument matched by varags)
            specifies the minimum time {\bf sck} or {\bf sdo} must remain
            unchanged following a transition for that transition to be
            recognised, i.e. it is a deglitch time period. The argument may be
            an immediate integer, in which case it specifies the number of clock
            cycles in the current clock domain, or it may be immediate floating
            point, which will be interpreted as a time in second. If the argument
            is missing or is zero no deglitching will be applied. Note that
            deglitching will introduce a data output latency.
            
            The clock domain within which this module is called must be
            of a frequency high enough to properly sample {\bf sck} and
            {\bf sdo}, i.e. at least twice the frequency of {\bf sck}.

            See also library module {\bf tws\_send()} \autorefph{sec:lmtwssend}.



        \subsubsection{tws\_send - two wire serial transmitter}\label{sec:lmtwssend}

            {\bf Library: stdlib.3pl}


            {\bf Prototype:} \\
            \vsp
            \begin{lstlisting}[gobble=12, language=threepl]
               tws_send(
                   static(sck, "log"), static(sdo, "log")
                   <-
                   queue(in), varargs
               ) {}
            \end{lstlisting}

            {\bf Output parameters:} \\
            {\bf sck} is the output serial clock. \\
            {\bf sdo} is the output serial data.

            {\bf Input parameters:} \\
            {\bf in} is the input queue variable. \\
            {\bf varargs} matches two required arguments, {\bf timebase} and
            {\bf silence}.

            {\bf description:} \\
            This module transmits a two-wire serial stream from
            a parallel input on a queue. Output {\bf sck} is the serial clock
            and output {\bf sdo} is the serial data. The data changes
            on the negative edges of the clock.{\bf sck} is low between
            words.

            Argument {\bf timebase} (first argument matched by varags)
            may be immediate, value or static mode.
            
            If {\bf timebase} is immediate mode it must be type "float", "int"
            or "uint". If type "float" it specifies the period of the output
            clock {\bf sck} in second. If it is type "int" or "uint" it gives
            the period of {\bf sck} as a multiple of the period of the
            current clock.
            
            If {\bf timebase} is value or static mode it must be a stream
            of pulses in the current clock domain, each pulse initiating
            the next sck transition during serial transmission. If
            {\bf timebase} is continuously high then {\bf sck} will be at
            half the current clock rate.
            
            Argument {\bf silence} specifies the minimum gap between transmitted
            words. The argument may be an immediate integer, in which case it
            specifies the number of clock cycles in the current clock domain, or
            it may be immediate floating point, which will be interpreted as a
            time in second. If the argument is zero no gap between successive words
            will be inserted.
            
            The clock domain within which this module is called must be
            at least twice the frequency of {\bf sck}.

            See also library module {\bf tws\_recv()} \autorefph{sec:lmtwsrecv}.


    \subsection{library procedures}
            





        \subsubsection{assign - generalised assignment}\label{sec:lpassign}

            {\bf Library: stdlib.3pl}

            {\bf Prototype:} \\
            \vsp
            \begin{lstlisting}[gobble=12, language=threepl]
               global procedure
               assign (var(out) <- varargs) {}
            \end{lstlisting}

            {\bf Output parameters:} \\
            {\bf out} is the variable to be assigned.

            {\bf Input parameters:} \\            
            {\bf varargs} expects a single argument which is the input value.
        
            {\bf In-line target execution:} \\
            yes

            {\bf description:} \\
            This procedure is provided for cases where assignment is to be
            made to a variable whose mode is not known or which varies from
            instance to instance.
            
            The output variable may be of mode immediate, value, priority,
            clock, static, queue or selectvalue. The input may be any expression,
            immediate or target mode and in the latter case may contain queue
            reads.





        \subsubsection{cinclude - convert a C include file}\label{sec:lpcinclude}

            {\bf Library: stdlib.3pl}

            {\bf Prototype:} \\
            \vsp
            \begin{lstlisting}[gobble=12, language=threepl]
               global procedure
               cinclude(str(filename)) {}
            \end{lstlisting}

            {\bf Output parameters:} \\
            none

            {\bf Input parameters:} \\            
            {\bf filename} is the C include file name.
        
            {\bf In-line target execution:} \\
            no

            {\bf description:} \\
            This procedure reads a C include file and defines variables
            corresponding to each \#define entry that it finds. This enables
            a C include file to be shared between a CPU C program and
            an associated 3PL FPGA program where there is a CPU/FPGA interface
            and constants such as bus or register addresses must be the same.

            Note that {\bf cinclude()} only handles \#define, C-comments and
            blank lines. Anything else gives a fatal error.
            
            {\bf cinclude()} processes \#define definitions with integer,
            string and floating point definitions. It also handles a
            \#define with no value, creating a {\bf log} type initialised
            to {\bf true}.
            
            {\bf cinclude()} ignores C style comments (which must not extend
            beyond one line), C++ style comments and empty lines.
            
            Lines beginning with \#define must contain an identifier and an
            optional associated constant. Integer constants may be decimal,
            hexadecimal (0x\dots or 0X\dots), octal (0\dots) or binary (0b\dots or
            0B\dots). String constants must be enclosed quotes -  "\dots".
            
            \#defines cannot be nested.
            
            Each definition produces a corresponding global immediate {\bf int},
            {\bf str}, {\bf float} or {\bf log} as appropriate whose initial value
            is the associated value from the definition. If there is no value a
            {\bf log} type is created initialised to {\bf true}.
            
            For a C include file def.h as follows -
            \begin{lstlisting}[gobble=12, language=C]
                   /* Include file */
                   /* has only single line C style comments */

                   #define VARA    3   // also ignores c++ style comments
                   #define VARB    0xf /* second */
                   #define VARC    "name"
                   #define VARD    5.67
                   #define ENABLE
            \end{lstlisting}

            The 3PL source
            \begin{lstlisting}[gobble=12, language=threepl]
               cinclude("def.h");
            \end{lstlisting}
            will read the file def.h and effectively execute the following
            in the global module -
            \begin{lstlisting}[gobble=12, language=threepl]
               int(VARA, 3);
               int(VARB, 0xf);
               str(VARC, "name");
               float(VARD, 5.67);
               log(ENABLE, true);
            \end{lstlisting}


        \subsubsection{netlistcommentfromdirs - directives to nelist comment}\label{sec:lpncfromd}

            {\bf Library: stdlib.3pl}

            {\bf Prototype:} \\
            \vsp
            \begin{lstlisting}[gobble=12, language=threepl]
               global procedure
               netlistcommentfromdirs () {}
            \end{lstlisting}

            {\bf Output parameters:} \\
            none.

            {\bf Input parameters:} \\            
            none.
        
            {\bf In-line target execution:} \\
            no

            {\bf description:} \\
            This procedure copies a listing of the current identifiers and values
            of all directives to the directive "netlistcomment". Directive
            "netlistcomment" if defined is copied as comments to the output netlist
            and constraint files. If this procedure is called at the beginning
            of source code only directives defined by -D command line options
            and the directive "sourceversion" will be listed. If called at the end
            of execution then directives "part" and "family" will be listed
            along with the optinal "author" directive and any other directives
            defined by the program.


        \subsubsection{queuearray - construct an array of queues}\label{sec:lpqueuearray}

            {\bf Library: stdlib.3pl}

            {\bf Prototype:} \\
            \vsp
            \begin{lstlisting}[gobble=12, language=threepl]
               global procedure
               queuearray(var(ptr) <- type(type)) {}
            \end{lstlisting}

            {\bf Output parameters:} \\
            {\bf ptr} is an array of pointers to queues. Each
            queue has the type given by the input parameter.

            {\bf Input parameters:} \\            
            {\bf type} is the type of the queues.
        
            {\bf In-line target execution:} \\
            no

            {\bf description:} \\
            Create an array of queues. Since 'queue' is a mode, not a type,
            there cannot be an array of them, hence an array of pointers to
            queues must be used instead (the queue type itself can be an
            array, but that is not what is wanted in this case).
            
            The pointers to created queues are assigned to the output
            argument which must be a pointer or pointer array. The array may
            have more than one dimension. The queue type is given by the input
            argument which may be a type or a string equivalent. The created
            queues have the default depth of 2. If the output argument is a
            single pointer then a single new queue will be assigned to that
            pointer.



        \subsubsection{prelementIntPropertys - print a list of variable attributes}\label{sec:lpprintattrs}

            {\bf Library: stdlib.3pl}

            {\bf Prototype:} \\
            \vsp
            \begin{lstlisting}[gobble=12, language=threepl]
               global procedure
               prelementIntPropertys (str(lv, "[](key, str, mode, str, value, str)")) {
            \end{lstlisting}

            {\bf Output parameters:} \\
            none

            {\bf Input parameters:} \\            
            A call to listattributes().
        
            {\bf In-line target execution:} \\
            no

            {\bf description:} \\
            This procedure prints a list of attributes for a variable. Each line
            list in order -
            \begin{lstlisting}[gobble=12, language=threepl]
               key mode value
            \end{lstlisting}
            the fields being separated by horizontal tabs. The format is
            not very sophisticated.
        
            The argument must be a call to inbuilt function
            {\bf listattributes()} \autorefph{sec:iflistattrs} which itself
            has one argument, the variable identifier.



        \subsubsection{printdirs - print a list of directives}\label{sec:lpprintdirs}

            {\bf Library: stdlib.3pl}

            {\bf Prototype:} \\
            \vsp
            \begin{lstlisting}[gobble=12, language=threepl]
               global procedure
               printdirs () {}
            \end{lstlisting}

            {\bf Output parameters:} \\
            none

            {\bf Input parameters:} \\            
            none.
        
            {\bf In-line target execution:} \\
            no

            {\bf description:} \\
            This procedure prints a list of directives. Each line list
            in order -
            \begin{lstlisting}[gobble=12, language=threepl]
               identifier type value
            \end{lstlisting}
            the fields being separated by horizontal tabs. The format is
            not very sophisticated.




        \subsubsection{printvars - print a list of variables}\label{sec:lpprintvars}

            {\bf Library: stdlib.3pl}

            {\bf Prototype:} \\
            \vsp
            \begin{lstlisting}[gobble=12, language=threepl]
               global procedure
               printvars (var(lv, "[](scope, str, ident, str, var, ->)")) {}
            \end{lstlisting}

            {\bf Output parameters:} \\
            none

            {\bf Input parameters:} \\            
            A call to listvars().
        
            {\bf In-line target execution:} \\
            no

            {\bf description:} \\
            This procedure prints a list of variables. Each line list
            in order -
            \begin{lstlisting}[gobble=12, language=threepl]
               scope identifier mode type
            \end{lstlisting}
            the fields being separated by horizontal tabs. The format is
            not very sophisticated.

            The argument must be a call to inbuilt function
            {\bf listvars()} \autorefph{sec:iflistvars} which itself
            has an optional argument.

            Argument "global" causes only global variables to be printed.

            Argument "file" causes only file scope variables to be printed.

            Argument "files" causes variables in all file scopes to be printed.

            Argument "local" causes only local variables to be printed, i.e. only
            those declared within the current module, procedure or function, and
            not variables in nested control structures.

            Argument "" selects the default, i.e. all block and local variables,
            the file scope variables and global scope variables will be printed.

            See also inbuilt function {\bf listvars()} \autorefph{sec:iflistvars}.

%
%        \subsubsection{rread2 - read from an extended dual-port registered
%                                            output memory}\label{sec:lprread2}
%
%            {\bf Library: stdlib.3pl}
%
%            {\bf Prototype:} \\
%            \vsp
%            \begin{lstlisting}[gobble=12, language=threepl]
%               global procedure
%               rread2 (var(m, "->"), int(port), value(addr)) {}
%            \end{lstlisting}
%
%            {\bf Output parameters:} \\
%            none
%
%            {\bf Input parameters:} \\            
%            {\bf m} is a pointer to the registered output memory struct.\\
%            {\bf port} is the port.\\
%            {\bf addr} is the read address.
%        
%            {\bf In-line target execution:} \\
%            yes
%
%            {\bf description:} \\
%            This procedure reads from an extended registered output memory that
%            has been created by the function {\em bram2()} \autorefph{sec:lfbram2}. The
%            read is not arbitrated and the application code must ensure that
%            reads and writes are mutually exclusive. There may be any number of
%            calls to this procedure and the associated write procedure {\em
%            rwrite2()} \autorefph{sec:lprwrite2} for the same port of the same memory as
%            long as their target execution is mutually exclusive.
%
%            The address expression must match the type given by the address type
%            argument to {\em bram2()} \autorefph{sec:lfbram2} and may be immediate, value,
%            selectvalue, static or queue mode. (for primitive types the bit width
%            may differ).
%            
%            The value resulting from the read will be available from the
%            following clock cycle at the output of the memory. This value can
%            be accessed via the library function {\em rramout2} \autorefph{sec:lfrramout2}.
%
%
%        \subsubsection{rwrite2 - write to an extended dual-port registered
%                                            output memory}\label{sec:lprwrite2}
%
%            {\bf Library: stdlib.3pl}
%
%            {\bf Prototype:} \\
%            \vsp
%            \begin{lstlisting}[gobble=12, language=threepl]
%               global procedure
%               rwrite2 (var(m, "->"), int(port), value(addr), value(data)) {}
%            \end{lstlisting}
%
%            {\bf Output parameters:} \\
%            none
%
%            {\bf Input parameters:} \\            
%            {\bf m} is a pointer to the registered output memory struct.\\
%            {\bf port} is the port.\\
%            {\bf addr} is the write address.\\
%            {\bf data} is the data to be written.
%        
%            {\bf In-line target execution:} \\
%            yes
%
%            {\bf description:} \\
%            This procedure writes to an extended registered output memory that
%            has been created by the function {\em bram2()} \autorefph{sec:lfbram2}. The
%            write is not arbitrated and the application code must ensure that
%            reads and writes are mutually exclusive. There may be any number of
%            calls to this procedure and the associated read procedure {\em
%            rread2()} \autorefph{sec:lprread2} for the same port of the same memory as
%            long as their target execution is mutually exclusive.
%            
%            The registered output of the port is also changed to the written
%            value.
%
%            The address and data expressions must match the types given by
%            the arguments to {\em bram2()} \autorefph{sec:lfbram2} and may be
%            immediate, value, selectvalue, static or queue mode (for primitive
%            types the bit width may differ).
%

        \subsubsection{wait - wait for a time period}\label{sec:lpwait}

            {\bf Library: stdlib.3pl}

            {\bf Prototype:} \\
            \vsp
            \begin{lstlisting}[gobble=12, language=threepl]
               global procedure
               wait (var(t)) {}
            \end{lstlisting}

            {\bf Output parameters:} \\
            none

            {\bf Input parameters:} \\            
            {\bf t} is the delay in cycles or second.
        
            {\bf In-line target execution:} \\
            yes

            {\bf description:} \\
            This procedure waits until a time has elapsed. The input
            argument may be immediate, value or static mode. If immediate
            the wait period is determined on immediate execution. If
            a target mode the wait period is given by the value of the
            argument at the time the procedure is executed in the FPGA.
            
            If the input argument is immediate mode, it must be type {\bf
            uint}, {\bf int} or {\bf float}. If the argument is type {\bf
            float} it is the number of second to wait. The delay will be
            rounded to the nearest integer number of clock cycles of the
            current clock. If the argument is an integer it is the number of
            clock cycles to wait.

            If the input argument is value or static mode it must be type {\bf
            uint} or {\bf int} and is the number of clock cycles to wait. The
            value must not be zero - this will generate the longest wait
            consistent with the width of the argument variable.
           


        \subsubsection{wait\_until - wait on a logical value}\label{sec:lpwaituntil}

            {\bf Library: stdlib.3pl}

            {\bf Prototype:} \\
            \vsp
            \begin{lstlisting}[gobble=12, language=threepl]
               global procedure
               wait_until (value(v, "log")) {}
            \end{lstlisting}

            {\bf Output parameters:} \\
            none

            {\bf Input parameters:} \\            
            {\bf v} is the input value to wait on.
        
            {\bf In-line target execution:} \\
            yes

            {\bf description:} \\
            This procedure waits until the logical input argument 'v' is
            true.


        \subsubsection{wait\_while - wait on a logical value}\label{sec:lpwaitwhile}

            {\bf Library: stdlib.3pl}

            {\bf Prototype:} \\
            \vsp
            \begin{lstlisting}[gobble=12, language=threepl]
               global procedure
               wait_while (value(v, "log")) {}
            \end{lstlisting}

            {\bf Output parameters:} \\
            none

            {\bf Input parameters:} \\            
            {\bf v} is the input value to wait on.
        
            {\bf In-line target execution:} \\
            yes

            {\bf description:} \\
            This procedure waits while the logical input argument 'v' is
            true.





    \subsection{library functions}




        \subsubsection{argstomap - create a map from key=value arguments}\label{sec:lfargstomap}

            {\bf Library: stdlib.3pl}

            {\bf Prototype:} \\
            \vsp
            \begin{lstlisting}[gobble=12, language=threepl]
               global function
               argstomap (varargs) {}
            \end{lstlisting}

            {\bf Return mode and type:} \\
            immediate, map

            {\bf Input parameters:} \\            
            See description.

            {\bf description:} \\
            This function takes a number of key=value arguments and from them
            builds a map.



 
            


        \subsubsection{asyncqueuespaces - space count for asynchronous queue}\label{sec:lfasyncqueuespaces}

            {\bf Library: stdlib.3pl}

            {\bf Prototype:} \\
            \vsp
            \begin{lstlisting}[gobble=12, language=threepl]
               global function
               asyncqueuespaces (queue(p)) {}
            \end{lstlisting}

            {\bf Return mode and type:} \\
            value, "uint:"

            {\bf Input parameters:} \\            
            {\bf p} is an asynchronous queue variable.

            {\bf description:} \\
            This function returns a target value which gives a guaranteed
            minimum number of free spaces in the asynchronous queue, rather than
            a precise count. The count is determined from status bits derived
            from the most significant two bits of the of the read and write gray
            coded addresses  (see inbuilt function {\bf queuewritestatus(})
            (\autorefph{sec:ifqueuewritestatus}). This uses much less logic than
            an exact determination of the word count.
            
            This function is intended for an external interface to the FPGA
            which writes a queue. The count can be used to avoid the extra
            transfers which would be required to determine non-full status
            prior to every write.
            
            This must be called in the write clock domain.
             


        \subsubsection{asyncqueuewords - word count for asynchronous queue}\label{sec:lfasyncqueuewords}

            {\bf Library: stdlib.3pl}

            {\bf Prototype:} \\
            \vsp
            \begin{lstlisting}[gobble=12, language=threepl]
               global function
               asyncqueuewords (queue(p)) {}
            \end{lstlisting}

            {\bf Return mode and type:} \\
            value, "uint:"

            {\bf Input parameters:} \\            
            {\bf p} is an asynchronous queue variable.

            {\bf description:} \\
            This function returns a target value which gives a guaranteed
            minimum number of words in the asynchronous queue, rather than a
            precise count. The count is determined from status bits derived from
            the most significant two bits of the of the read and write gray
            coded addresses (see inbuilt function {\bf queuereadstatus(}
            (\autorefph{sec:ifqueuereadstatus})). This uses much less logic than
            an exact determination of the word count.
            
            This function is intended for an external interface to the FPGA
            which reads a queue. The count can be used to avoid the extra
            transfers which would be required to determine non-empty status
            prior to every read.
            
            This must be called in the read clock domain.
           


        \subsubsection{bin\_to\_gray - convert binary count to gray count}\label{sec:lfbintogray}

            {\bf Library: stdlib.3pl}

            {\bf Prototype:} \\
            \vsp
            \begin{lstlisting}[gobble=12, language=threepl]
               global function
               bin_to_gray (value(b)) {}
            \end{lstlisting}

            {\bf Return mode and type:} \\
            immediate, \verb'->'

            {\bf Input parameters:} \\            
            {\bf b} is the input binary count value.

            {\bf description:} \\
            This function converts an input weighted binary value to




        \subsubsection{bitset - test if a bit is set}\label{sec:lfbitset}

            {\bf Library: stdlib.3pl}

            {\bf Prototype:} \\
            \vsp
            \begin{lstlisting}[gobble=12, language=threepl]
               global function
               bitset (value(in), int(bit)) {}
            \end{lstlisting}

            {\bf Input parameters:} \\            
            {\bf in} is the input value. It may be any target type.\\
            {\bf bit} is the index of the required bit.

            {\bf return mode:} \\
            value

            {\bf return type:} \\
            log

            {\bf Description:} \\
            Returns true if a selected bit from the first (target) argument is
            set. The second argument is an immediate integer giving the bit
            index, 0 selecting the LSB. The first argument type need not be
            primitive, but it is better coding practice to select a bit from an
            array member or struct field rather than from the compound type.
            
            See also {\bf getbit()} \autorefph{sec:lfgetbit} which does the same
            but returns as type "bits:1".
            





        \subsubsection{deglitch - deglitch an input signal}\label{sec:lfdeglitch}

            {\bf Library: stdlib.3pl}

            {\bf Prototype:} \\
            \vsp
            \begin{lstlisting}[gobble=12, language=threepl]
               global function
               deglitch (value(in), varargs) {}
            \end{lstlisting}

            {\bf Return mode and type:} \\
            value, same type as input

            {\bf Input parameters:} \\            
            {\bf in} is the input signal. \\
            {\bf varargs} expects a single argument (at most) giving the
            deglitch time constant.
            
            The output of this function follows the input value, however
            the output will not follow an input change until the input
            has maintained its new value for a time specified by the
            second argument.
            
            The first argument is the input signal or value. It may be any
            target type. It would normally be input mode or an expression
            containing an input mode variable.
            
            Other modes are allowed for the input (except cmemory, rmemory or
            clock mode) and it may contain queue reads, but there would not seem
            to be any reason to use any of these.
            
            The second argument may be missing, in which case the output is just
            connected to the input and there is no deglitching performed. The
            second argument may be immediate mode of type "int", "uint" or
            "float", or static mode of type "int:" or "uint:". If integer it
            specifies the deglitching time in terms of clock cycles of the
            current clock. If it is floating point the argument specifies an
            explicit time in second and the delay will be of the minimum
            number of clock cycles which encompass the delay time. The second
            argument will be ignored if it is immediate and the number of clock
            cycles specified, or the number derived from the floating point
            time, is 1 or 0 (or -ve). In that case the output is again simply
            connected to the input.




        \subsubsection{delay - fixed delay}\label{sec:lfdelay}

            {\bf Library: stdlib.3pl}

            {\bf Prototype:} \\
            \vsp
            \begin{lstlisting}[gobble=12, language=threepl]
               global function
               delay (value(in), int(len), value(step, "log"), log(ram)) {}
            \end{lstlisting}

            {\bf Return mode and type:} \\
            value, same type as input

            {\bf Input parameters:} \\            
            {\bf in} is the input signal. \\
            {\bf len} is the length of the delay line.\\
            {\bf step} is an optional step control.\\
            {\bf ram} is an optional switch to force usage of RAM

            This has up to 4 input arguments. Input argument {\bf len} is
            immediate mode. {\bf step} may be immediate or target mode as
            described below. The data input argument {\bf in} must not be modes
            {\bf immediate}, {\bf queue}, {\bf memory} or {\bf clock} and value
            mode arguments must not contain queue reads.
            
            The output of this function follows the input value delayed by {\bf
            del} clock cycles.
            
            The argument to {\bf step} must be type "log" and may be as follows -
            \blist
            \item   no argument - the delay will advance on each clock cycle
            \item   an immediate mode value (or constant value mode) which is
                    {\bf true} - the delay will advance on each clock cycle
            \item   an immediate mode value (or constant value mode) which is
                    {\bf false} - the delay will advance on clock cycles in
                    which the function is evaluated in the target code, e.g.
                    on used in assignment
            \item   a static, selectvalue or value mode value - the delay will
                    advance on clock cycles when {\bf step} is true.
            \blend
            For example -
            \begin{lstlisting}[gobble=12, language=threepl]
                when (l)
                    a <- delay(b, 10, false);
            \end{lstlisting}
            will shift the current value of {\bf b} into the delay line every
            time the assignment to {\bf a} is made ({\bf l} is true), the
            assigned value being popped off the end.
           
            {\bf step} false can only be used where {\bf delay()} is evaluated
            in a statement executing within an execution thread. For example -
            \begin{lstlisting}[gobble=12, language=threepl]
                output(,, delay(b, 10, false), loc=A_pins[0]);
            \end{lstlisting}
            would give an error (function used() called in illegal context).
            
            Whereas -
            \begin{lstlisting}[gobble=12, language=threepl]
                when (l)
                    a <- delay(b, 10);
            \end{lstlisting}
            will shift {\bf b} into the delay line every clock cycle, popping
            off the vaue at the other end. When {\bf l} is true the value popped
            off will also be assigned to {\bf a}.
            \begin{lstlisting}[gobble=12, language=threepl]
                when (l)
                    a <- delay(b, 10, true);
            \end{lstlisting}
            will behave the same way.

            Where {\bf step} is variable during target execution the shifting
            is controlled solely by step while the assignment, controlled by
            {\bf l}, assigns the end value of the delay line to {\bf a}
            regardless of whether a shift has occured or not.
            \begin{lstlisting}[gobble=12, language=threepl]
                when (l)
                    a <- delay(b, 10, step);
            \end{lstlisting}
            
            The data width is not limited. This function can operate as a
            sequence generator if the output is connected back to the input.
            
            For the first {\bf del} steps the output will be 0, false etc.
            
            {\bf in} may be any type. It must not contain queue reads (fatal
            error).
            
            The return value is the same type as {\bf in} and is the delayed
            input.
            
            If the delay is small enough (see table below) this will call
            inbuilt module {\bf del()} \autorefph{sec:imdel} to implement the
            delay using CLB shift registers, otherwise rmemory (block RAM) will
            be used. In the latter case the length {\bf len} is limited to the
            maximum allowed memory depth for the FPGA family -
            
            \begin{tabular}{| l | l | l |} \hline
                                     &  CLB SR maximum length & maximum RAM depth             \\ \hline \hline
            Spartan-II, Spartan IIE,
            Virtex, Virtex-E         &   16 & 4096 \\ \hline
            Spartan-3, Virtex-II,
            Virtex-II Pro,
            Virtex 4,                &  16 & 16384 \\ \hline
            Spartan 6, Virtex 5,
            Virtex 6, series 7, Zynq &  32 & 32768 \\ \hline
            \end{tabular}

            If parameter {\bf ram} is true block RAM will be used for the
            delay despite a small delay being specified.


            See also library function {\bf del()} \autorefph{sec:lfdel} in
            library xilinx/common which allows a variable-length delay line but
            does not offer the option of using block RAM rather than shift
            registers. That section also ilustrates a way of using a queue input,
            getting around the no-queue limitation described above.


        \subsubsection{forcecast - force a prohibited cast}\label{sec:lfforcecast}

            {\bf Library: stdlib.3pl}

            {\bf Prototype:} \\
            \vsp
            \begin{lstlisting}[gobble=12, language=threepl]
               global function
               forcecast (value(v), type(t)) {}
             \end{lstlisting}

            {\bf Input parameters:} \\            
            {\bf v} is the input value. It may be any target type.\\
            {\bf t} is the target type to which the input value is to be cast.

            {\bf return mode:} \\
            value

            {\bf return type:} \\
            given by 2nd argument

            {\bf Description:} \\
            Returns the target first argument value cast to the target type
            given by the second argument.
            
            This function can be used to perform an otherwise prohibited
            cast operation. It does this by casting in two steps via an
            intermediate "bits:" type. This should be used with caution!




        \subsubsection{getbit - get a bit}\label{sec:lfgetbit}

            {\bf Library: stdlib.3pl}

            {\bf Prototype:} \\
            \vsp
            \begin{lstlisting}[gobble=12, language=threepl]
               global function
               getbit (value(in), int(bit)) {}
            \end{lstlisting}

            {\bf Input parameters:} \\            
            {\bf in} is the input value. It may be any target type.\\
            {\bf bit} is the index of the required bit.

            {\bf return mode:} \\
            value

            {\bf return type:} \\
            bits:1

            {\bf Description:} \\
            Returns a selected bit from the first (target) argument. The second
            argument is an immediate integer giving the bit index, 0 selecting
            the LSB. The first argument type need not be primitive, but
            it is better coding practice to select a bit from an array member
            or struct field rather than from the compound type.
            
            See also {\bf bitset()} \autorefph{sec:lfbitset} which does the same
            but returns as type "log".
    



        \subsubsection{incrdecr - dynamically increment or decrement}\label{sec:lfincrdecr}

            {\bf Library: stdlib.3pl}

            {\bf Prototype:} \\
            \vsp
            \begin{lstlisting}[gobble=12, language=threepl]
               global function
               incrdecr (value(in), value(incr, "log")) {}
            \end{lstlisting}

            {\bf Return mode and type:} \\
            value, type(in)

            {\bf Input parameters:} \\            
            {\bf in} is the input value to increment or decrement. It may be
            types int or uint.\\
            {\bf incr} selects incrementing (true) or decrementing (false).

            {\bf description:} \\
            Increment or decrement a value according to a logical input.



        \subsubsection{isclear - verify a run of clear bits}\label{sec:lfisclear}

            {\bf Library: stdlib.3pl}

            {\bf Prototype:} \\
            \vsp
            \begin{lstlisting}[gobble=12, language=threepl]
               global function
               isclear (value(in), int(bit), int(len, 1)) {}
            \end{lstlisting}

            {\bf Return mode and type:} \\
            value, log

            {\bf Input parameters:} \\            
            {\bf in} is the value to be tested. \\
            {\bf bit} is the bit index from the least significant end. \\
            {\bf len} is the number of bits.

            {\bf description:} \\
            This target function returns true if {\bf len} bits starting at
            bit {\bf bit} are all clear.



        \subsubsection{isset - verify a run of set bits}\label{sec:lfisset}

            {\bf Library: stdlib.3pl}

            {\bf Prototype:} \\
            \vsp
            \begin{lstlisting}[gobble=12, language=threepl]
               global function
               isset (value(in), int(bit), int(len, 1)) {}
            \end{lstlisting}

            {\bf Return mode and type:} \\
            value, log

            {\bf Input parameters:} \\            
            {\bf in} is the value to be tested. \\
            {\bf bit} is the bit index from the least significant end. \\
            {\bf len} is the number of bits.

            {\bf description:} \\
            This target function returns true if {\bf len} bits starting at
            bit {\bf bit} are all set.



        \subsubsection{lfsr - linear feedback shift register}\label{sec:lfsr}

            {\bf Library: stdlib.3pl}

            {\bf Prototype:} \\
            \vsp
            \begin{lstlisting}[gobble=12, language=threepl]
               global function
               lfsr (var(t), value(step, "log", true), int(init)) {}
            \end{lstlisting}

            {\bf Return mode and type:} \\
            value, uint

            {\bf Input parameters:} \\            
            {\bf t} is the result type. \\
            {\bf step} is an optional control input to advance the output. \\
            {\bf init} is an optional initial value.

            {\bf description:} \\
            This target function returns a pseudo-random number. The first
            argument specifies the type of the function value. It may be
            a type "str" or type "type". If the former it must be a valid
            type specification string. Any output type may be specified though
            it is difficult to imagine occasions where a compound type would
            be appropriate. The size of the type must be in the range 3 to 32
            bits.
            
            If the second argument is provided it controls the advance of the
            pseudo-random sequence from one to the next. If absent the sequence
            will advance on every clock cycle.
            
            If the last parameter has an argument it provides an initial
            value for the register. This may have any value. If the value
            is larger than can be accommodated in the register it will be
            truncated at the most significant end. If the last parameter
            has no argument then the initial value of the register will
            be a random number derived from calls to inbuilt function
            random() for each bit.
            
            For an output type of size n bits all $2^n$ values are produced (there
            is no lock-up state, or rather the 'lock-up' state is explicitly forced
            and exited.



        \subsubsection{lszeros - count least significant zero bits}\label{sec:lflszeros}

            {\bf Library: stdlib.3pl}

            {\bf Prototype:} \\
            \vsp
            \begin{lstlisting}[gobble=12, language=threepl]
               global function
               lszeros (value(in, "uint:")) {}
            \end{lstlisting}

            {\bf Return mode and type:} \\
            value, uint

            {\bf Input parameters:} \\            
            {\bf in} is the operand.

            {\bf description:} \\
            This target function returns the number of least significant zero
            bits in a non-zero unsigned integer. If the argument is zero the
            result is also zero. This function was written for use in floating
            point number normalisation (the mantissa is the argument) but may
            find other uses.




        \subsubsection{maxval - target maximum value}\label{sec:lfmaxval}

            {\bf Library: stdlib.3pl}

            {\bf Prototype:} \\
            \vsp
            \begin{lstlisting}[gobble=12, language=threepl]
               global function
               maxval (varargs) {}
            \end{lstlisting}

            {\bf return mode:} \\
            immediate or value

            {\bf return type:} \\
            type of arguments

            {\bf Description:} \\
            Returns the maximum value of any number of immediate or target
            variables or expressions. The arguments must be type int or uint and
            must all be the same type. A mixture of immediate and target
            arguments is allowed, but in that case only a single immediate
            argument should be used as the function will not optimise the
            immediate arguments to a single immediate value.            


        \subsubsection{minval - target maximum value}\label{sec:lfminval}

            {\bf Library: stdlib.3pl}

            {\bf Prototype:} \\
            \vsp
            \begin{lstlisting}[gobble=12, language=threepl]
               global function
               minval (varargs) {}
            \end{lstlisting}

            {\bf return mode:} \\
            immediate or value

            {\bf return type:} \\
            type of arguments

            {\bf Description:} \\
            Returns the minimum value of any number of immediate or target
            variables or expressions. The arguments must be type int or uint and
            must all be the same type. A mixture of immediate and target
            arguments is allowed, but in that case only a single immediate
            argument should be used as the function will not optimise the
            immediate arguments to a single immediate value.            



        \subsubsection{multiplex\_struct\_l - multiplex output type with log array\\
                    multiplex\_struct\_e - multiplex output type with uint}\label{sec:lfmultiplexstruct}

            {\bf Library: stdlib.3pl}

            {\bf Prototypes:} \\
            \vsp
            \begin{lstlisting}[gobble=12, language=threepl]
               multiplex_struct_l (type(intype), int(n)) {}

               multiplex_struct_p (type(intype), int(n)) {}
            \end{lstlisting}

            {\bf return mode:} \\
            immediate

            {\bf return type:} \\
            a struct type

            {\bf Description:} \\
            These functions are used in conjunction with library modules
            {\bf multiplex()} \autorefph{sec:lmmultiplex} \\
            {\bf multiplex\_l()} \autorefph{sec:lmmultiplex} \\
            {\bf multiplex\_e()} \autorefph{sec:lmmultiplex} \\
            {\bf multiplexp()} \autorefph{sec:lmmultiplex} \\
            {\bf multiplexp\_l()} \autorefph{sec:lmmultiplex} \\
            {\bf multiplexp\_e()} \autorefph{sec:lmmultiplex} \\
            {\bf multiplext()} \autorefph{sec:lmmultiplex} \\
            {\bf multiplext\_l()} \autorefph{sec:lmmultiplex} \\
            {\bf multiplext\_e()} \autorefph{sec:lmmultiplex} \\
            {\bf demultiplex()} \autorefph{sec:lmdemultiplex} \\
            {\bf demultiplex\_l()} \autorefph{sec:lmdemultiplex} \\
            {\bf demultiplex\_e()} \autorefph{sec:lmdemultiplex}
            They return a struct which is a
            wrapper for the output data type. The wrapper allows for an
            associated identifier for the source port of a datum. The struct has
            two fields. Field {\bf body} is the same type as the multiplexer
            inputs and contains a multiplexed datum. Field {\bf port} contains
            an identifier for the source port from which the datum originated.

            The first argument is the type of the inputs to module {\bf
            multiplex()} or {\bf multiplexp()} or the outputs of {\bf
            demultiplex()}.

            The second argument is the number of inputs to module {\bf
            multiplex()} or {\bf multiplexp()} or outputs to {\bf
            demultiplex()}.

            In the case of {\bf multiplex\_struct\_l()} field {\bf port} will be
            an array of type log whose dimension is equal to the number of
            inputs to the associated multiplexer or the outputs of the
            associated demultiplexer. This field has one member of the array
            true to identify the associated datum source port.

            In the case of {\bf multiplex\_struct\_e()} field {\bf port} will be
            a "uint" large enough to contain the port number (starting at 0).

            See also - \\
            {\bf multiplex()} \autorefph{sec:lmmultiplex} \\
            {\bf multiplex\_l()} \autorefph{sec:lmmultiplex} \\
            {\bf multiplex\_e()} \autorefph{sec:lmmultiplex} \\
            {\bf multiplexp()} \autorefph{sec:lmmultiplex} \\
            {\bf multiplexp\_l()} \autorefph{sec:lmmultiplex} \\
            {\bf multiplexp\_e()} \autorefph{sec:lmmultiplex} \\
            {\bf multiplext()} \autorefph{sec:lmmultiplex} \\
            {\bf multiplext\_l()} \autorefph{sec:lmmultiplex} \\
            {\bf multiplext\_e()} \autorefph{sec:lmmultiplex} \\
            {\bf demultiplex()} \autorefph{sec:lmdemultiplex} \\
            {\bf demultiplex\_l()} \autorefph{sec:lmdemultiplex} \\
            {\bf demultiplex\_e()} \autorefph{sec:lmdemultiplex}


        \subsubsection{mux - multiplex inputs}\label{sec:lfmux}

            {\bf Library: stdlib.3pl}

            {\bf Prototype:} \\
            \vsp
            \begin{lstlisting}[gobble=12, language=threepl]
               global function
               mux (value(s), varargs) {}
            \end{lstlisting}

            {\bf return mode:} \\
            value

            {\bf return type:} \\
            type of data argument

            {\bf Description:} \\
            Returns one of the data inputs. The first argument must be
            target mode of type {\bf uint} and selects which of the following
            input arguments is to be returned as the value of the function.
            The remaining arguments are the data inputs.

            If the data input expressions are primitive they must be the same
            type but their numeric widths may differ. The function value will be
            the width of the widest data input. If the data inputs are not
            primitive then they must all be exactly the same type, including
            primitive component widths.

        \subsubsection{ptr2expr - a pointer to an expression}\label{sec:lfptr2expr}

            {\bf Library: stdlib.3pl}

            {\bf Prototype:} \\
            \vsp
            \begin{lstlisting}[gobble=12, language=threepl]
               global function
               ptr2expr (varargs) {}
            \end{lstlisting}

            {\bf return mode:} \\
            value

            {\bf return type:} \\
            type of argument

            {\bf Description:} \\
            Returns a pointer to a value mode variable which has been
            assigned the target expression passed as the argument. The
            argument may also be a variable or constant. The argument
            cannot be mode clock or mode memory.


        \subsubsection{resync - resynchronise a pulse or level to another clock}\label{sec:lfresync}

            {\bf Library: stdlib.3pl}

            {\bf Prototype:} \\
            \vsp
            \begin{lstlisting}[gobble=12, language=threepl]
               global function
               resync (value(in, "log")) {}
            \end{lstlisting}

            {\bf return mode:} \\
            value

            {\bf return type:} \\
            log

            {\bf Description:} \\
            {\bf in} is an expression or a value, selectvalue, static or
            priority variable of type log. The argument must not contain queues.
        
            When the input is clocked true in its clock domain the function
            returned value will go true at the next positive edge of the
            current clock and will then go false at the following edge. This
            process cannot repeat until the input has returned to false for at
            least one clock cycle in its clock domain after the output has
            gone false.
            
            This function is used to produce a single cycle pulse in the
            current clock domain from a pulse or level in another clock
            domain.
            

%
%        \subsubsection{rramout2 - get the value of an extended dual-port
%                                registered output memory}\label{sec:lfrramout2}
%
%            {\bf Library: stdlib.3pl}
%
%            {\bf Prototype:} \\
%            \vsp
%            \begin{lstlisting}[gobble=12, language=threepl]
%               global function
%               rramout2 (var(m, "->"), int(port)) {}
%            \end{lstlisting}
%
%            {\bf Return mode and type:} \\
%            value, data type
%
%            {\bf Input parameters:} \\            
%            {\bf m} is a pointer to the registered output memory struct.\\
%            {\bf port} is the port.
%
%            {\bf description:} \\
%            This function returns a value which is the output register
%            of the memory port. The type returned is the data type passed
%            to {\em bram2()} \autorefph{sec:lfbram2} when creating the memory structure.
%





        \subsection{select - a combinatorial selector}\label{sec:lfselect}

            {\bf Library: stdlib.3pl}

            {\bf Prototype:} \\
            \vsp
            \begin{lstlisting}[gobble=12, language=threepl]
               global function
               select(s1, d1, s2, d2,...) {}
            \end{lstlisting}

            {\bf return mode:} \\
            value

            {\bf return type:} \\
            type of data inputs

            {\bf Description:} \\
            An inbuilt selector function using gate logic.
            This has any number of argument pairs.

            The 1st input argument of each pair is a logical expression which
            selects the following data input expression if true.

            The 2nd input argument of each pair is a data input signal. If the
            data input expressions are primitive they must be the same type but
            their numeric widths may differ. The function value will be the
            width of the widest data input. If the data inputs are not
            primitive then they must all be exactly the same type, including
            primitive component widths.

            The function has a defined value when at most one select input
            is available and true and the associated data input is available,
            the value reflecting the selected data input. If more than one
            input is selected the value is undefined. If no inputs are
            selected, the value is zero or false depending on data type(s).

            If none of the function inputs contain queues, the function value is
            always available. If any function inputs contain queues the function
            value is available when an input selector and its associated data
            input are available and the selector is true. Note that if some
            inputs contain queues but one input selector and its associated data
            do not, the function will be available when that selector is true.
            Note that the use of this function as a gate with a clear value
            when no selectors are true is only possible if the inputs contain
            no queues since any queues in any inputs will render the function
            value unavailable when no input is selected.
            
            See also inbuilt function {\bf selecta()} \autorefph{sec:ifselecta}
            where the select and data inputs are arrays of values or pointers.




        \subsubsection{sprintf - format a string from variables}\label{sec:lfsprintf}

            {\bf Library: stdlib.3pl}

            {\bf Prototype:} \\
            \vsp
            \begin{lstlisting}[gobble=12, language=threepl]
               global function
               sprintf (fmt, v1, v2, ...) {}
            \end{lstlisting}

            {\bf return mode:} \\
            immediate

            {\bf return type:} \\
            str

            {\bf Description:} \\
            This function is a wrapper for the inbuilt procedure
            {\bf sprintf()} \autorefph{sec:ipsprintf}. It returns the formatted string.
        
            The format argument may be followed by a single argument which is an array of
            pointers, for example {\bf inargs} in the calling module, procedure or function.
            In that case the pointers are resolved and provide the input variable list to be
            printed.







%        \subsubsection{svn\_version - get the current subversion version number}\label{sec:lfsvnversion}
%
%            {\bf Library: stdlib.3pl}
%
%            {\bf Prototype:} \\
%            \vsp
%            \begin{lstlisting}[gobble=12, language=threepl]
%               global function
%               svn_version () {}
%            \end{lstlisting}
%
%            {\bf Return mode and type:} \\
%            immediate, "str"
%
%            {\bf Input parameters:} \\            
%            None.
%
%            {\bf description:} \\
%            This function returns the current subversion version number as a
%            string. It does this by running the command "svnversion" in a shell in
%            the current directory and returning the resulting string from the shell
%            standard output with the trailing newline removed.


        \subsubsection{to\_pulse - generate a pulse from a transition}\label{sec:lftopulse}

            {\bf Library: stdlib.3pl}

            {\bf Prototype:} \\
            \vsp
            \begin{lstlisting}[gobble=12, language=threepl]
               global function
               to_pulse (value(in, "log")) {}
            \end{lstlisting}

            {\bf Return mode and type:} \\
            value, "log"

            {\bf Input parameters:} \\
            {\bf in} is the input value.

            {\bf description:} \\
            Generate a pulse out when the input argument {\bf in}
            changes from false to true.


        \subsubsection{truthtable - an expression from a truth table}\label{sec:lftruthtable}

            {\bf Library: stdlib.3pl}

            {\bf Prototype:} \\
            \vsp
            \begin{lstlisting}[gobble=12, language=threepl]
               global function
               truthtable (
                   var(table, "[64][6]str"),
                   value({i1, i2, i3, i4, i5, i6}, "log")
               ) {}
            \end{lstlisting}

            {\bf Return mode and type:} \\
            value, "log"

            {\bf Input parameters:} \\
            {\bf table} is the input value.
            {\bf i1} to {\bf i6} are the input values.

            {\bf description:} \\
            The first argument is the truth table as a 2-dimensional
            array of strings. This argument would usually be a 2-dimensional
            compound value, in which case its dimensions need not match the
            parameter dimensions as long as they do not exceed them. If the
            argument is a string array its dimensions must match those of the
            parameter, but the array would normally be initialised from
            a 2-dimensional compound value which again may have smaller
            dimensions. The string array dimensions are such as to allow
            a maximum of all 64 combinations of 6 variables.
            
            Truth table entries must be "0" or "1" for explicit entries
            or "", "x" or "X" for 'don`t care' entries. Compound values
            for string arrays may contain constant values of other types
            in which case these will be converted to strings. Thus a
            true entry may be given as {\bf "1"}, {\bf 1} or {\bf true},
            a false entry as {\bf "0"}, {\bf 0} or {\bf false} and a 'don`t
            care' entry may be simply omitted.
            
            The 2nd and following arguments are the input variables to the
            function described by the truth table. One to six of these
            arguments may be provided, trailing arguments being omitted.
            The arguments provided must be contiguous from {\bf i1}, i.e.
            there must be no omitted arguments except for trailing arguments.
            The number of arguments is used to determine how many columns of the
            truth table to process but no checks are made, e.g. extra columns
            are ignored and omitted columns will be treated as 'don`t care'
            entries.
            
            Example 1 - 4 input variables, compound value argument -
            \begin{lstlisting}[gobble=12, language=threepl]
               value({i, j, k, l, v}, "log");
               v := truthtable(
                       {  {1,  , 0, 1},
                          {0, 0, 1,  }  },
                           i, j, k, l
                    );
            \end{lstlisting}
            Note that the simplest and neatest truth table scheme is to
            use integers 0 and 1 for explicit entries (these will be converted
            to strings "0" and "1") and omit 'don`t care' entries. Note also
            that the input variables occur in the same order as the truth table
            columns and hence readability is aided by positioning them under
            the truth table when formatting the source code.
            
            Example 2 - 3 input variables, string array argument -
            \begin{lstlisting}[gobble=12, language=threepl]
               value({i, j, k, v}, "log");
               var(tt, "[64][6]str", {  {1,  , 0},
                                        {0, 0, 1}  });
               v := truthtable(tt, i, j, k);
            \end{lstlisting}
            Here a string array of the correct dimensions has been supplied
            as the truth table argument but the array initialisation is the
            same compound value as before. Of course the truth table string
            array could be generated by an algorithm if that were required.
            
            Example 3 - truth table function as a value variable initialiser -
            \begin{lstlisting}[gobble=12, language=threepl]
               value({i, j, k, l}, "log");
               value({v, "log", truthtable( {  {1,  , 0, 1},
                                               {0, 0, 1,  }  },
                                                i, j, k, l         );
            \end{lstlisting}



    \subsection{library classes}
    
        None.
